\documentclass[11pt]{article}

\usepackage[utf8]{inputenc}
\usepackage[T1]{fontenc}
\usepackage{geometry}
\usepackage{amsmath}
\usepackage{amssymb}
\usepackage{booktabs}
\usepackage{hyperref}
\usepackage{xurl}
\usepackage{xcolor}
\usepackage{listings}
\usepackage{enumitem}
\usepackage{parskip}
\usepackage{graphicx}
\graphicspath{{../images/}}

\geometry{margin=1in}

\lstset{
  language=Java,
  basicstyle=\ttfamily\small,
  keywordstyle=\color{blue},
  commentstyle=\color{gray},
  stringstyle=\color{teal},
  showstringspaces=false,
  columns=fullflexible,
  frame=single,
  framerule=0.2pt,
  breaklines=true
}

\setlist[itemize]{topsep=0.3em,itemsep=0.2em}
\setlist[enumerate]{topsep=0.3em,itemsep=0.2em}

% Simple per-section source line.
\newcommand{\notesources}[1]{\par\noindent{\small\color{gray}\textit{Sources:} \sloppy #1\par}}

% Unit-wise source strings for this subject (used by slides/notes).
% Path is relative to the lecture `latex/` folder (the compilation CWD).
% OOPS with Java (CSEG1044) — unit-wise sources
%
% These are short, student-facing reference strings intended to be used with:
%   - \slidesources{...}  (in beamer slides)
%   - \notesources{...}   (in lecture notes)
%
% Style inspiration: other subjects in My-Subjects/ use semicolon-separated sources
% and include an "accessed" date for web documentation.

\newcommand{\OOPSUnitISources}{Schildt, \textit{Java: A Beginner's Guide} (10e), McGraw-Hill (Java fundamentals + OOP basics).; Oracle Java Tutorials: Getting Started + Language Basics + Classes and Objects (accessed \today).; Java SE API docs (e.g., \texttt{String}, \texttt{Scanner}) (accessed \today).}

\newcommand{\OOPSUnitIISources}{Schildt, \textit{Java: The Complete Reference} (12e), McGraw-Hill (classes, inheritance, packages, interfaces).; Bloch, \textit{Effective Java} (3e), Addison-Wesley, 2018 (design choices; interfaces vs abstract classes).; Oracle Java Tutorials: Classes and Objects + Interfaces + Packages (accessed \today).; Java SE API docs (e.g., \texttt{Comparable}, \texttt{Runnable}) (accessed \today).}

\newcommand{\OOPSUnitIIISources}{Schildt, \textit{Java: The Complete Reference} (12e) (nested/inner classes, exceptions, threads, I/O).; Oracle Java Tutorials: Nested Classes + Exceptions + Concurrency + Basic I/O (accessed \today).; Java SE API docs (e.g., \texttt{Thread}, \texttt{AutoCloseable}) (accessed \today).}

\newcommand{\OOPSUnitIVSources}{Schildt, \textit{Java: The Complete Reference} (12e) (generics, lambdas, Swing, JDBC).; Bloch, \textit{Effective Java} (3e) (generics + lambdas).; Oracle Java Tutorials: Generics + Lambda Expressions + Swing + JDBC Basics (accessed \today).; Java SE API docs (e.g., \texttt{java.util.function}, \texttt{javax.swing}, \texttt{java.sql}) (accessed \today).}

\newcommand{\OOPSUnitVSources}{Schildt, \textit{Java: The Complete Reference} (12e) (Collections Framework + wrapper classes).; Oracle Java docs: Collections Framework overview (accessed \today).; Java SE API docs (\texttt{java.util} collections; wrapper classes in \texttt{java.lang}) (accessed \today).}

\newcommand{\OOPSUnitVISources}{Git SCM: \textit{Pro Git} (2e) (version control workflow), accessed \today.; GitHub Docs (repositories, issues, pull requests), accessed \today.; Oracle Java Tutorials: Swing + JDBC (accessed \today).; JUnit 5 User Guide (accessed \today).; Apache Maven Project: Getting Started (accessed \today).; Gradle User Manual: Getting Started (accessed \today).; UML class diagrams: UML 2.x overview/spec (accessed \today).}

\newcommand{\OOPSMiscWhyInterfacesSources}{Schildt, \textit{Java: The Complete Reference} (12e) (interfaces).; Bloch, \textit{Effective Java} (3e) (prefer interfaces where appropriate).; Oracle Java Tutorials: Interfaces (accessed \today).; Java SE API docs (e.g., \texttt{Comparable}, \texttt{Runnable}) (accessed \today).}



\title{Object Oriented Programming (CSEG1044)\\Unit 03 -- Lecture 02 Notes\\Exception Handling (Core)}
\begin{document}
\maketitle
\tableofcontents

\section{Lecture Snapshot}
\begin{itemize}
  \item \textbf{Unit focus:} Nested Classes, Exceptions, Multithreading, and I/O
  \item \textbf{Main new ideas:} Exceptions vs Errors, What happens when an exception occurs? (stack unwinding), How to read a stack trace (very important skill)
  \item \textbf{Course thread:} Use Exception Handling (Core) to make Campus Resource Manager safer, more modular, and easier to recover when something goes wrong.
\end{itemize}
\notesources{\OOPSUnitIIISources}

\section{Prerequisites}
You should already know:
\begin{itemize}
  \item methods, parameters, return values,
  \item classes and objects,
  \item basic input validation (e.g., checking empty strings).
\end{itemize}

\section{Learning Outcomes}
After this lecture, you should be able to:
\begin{itemize}
  \item Distinguish exception vs error; checked vs unchecked.
  \item Write \texttt{try/catch/finally} and use \texttt{throw} vs \texttt{throws}.
  \item Design a small custom exception for validation.
\end{itemize}

\section{Suggested Reading Order}
\begin{enumerate}
  \item Read the \textbf{Prerequisite Bridge} first and confirm each recalled idea.
  \item Study the central explanation in this order: \textit{Exceptions vs Errors} $\rightarrow$ \textit{What happens when an exception occurs? (stack unwinding)} $\rightarrow$ \textit{How to read a stack trace (very important skill)} $\rightarrow$ \textit{Checked vs unchecked exceptions (when to use which?)} $\rightarrow$ \textit{\texttt{try/catch/finally} and multiple catch}.
  \item Trace the worked example line by line before checking short answers.
  \item Attempt the practice section without looking at the solution immediately.
  \item Finish with the exit questions and further-study suggestions to confirm retention.
\end{enumerate}
\notesources{\OOPSUnitIIISources}

\section{Lecture Overview}
Exceptions are Java's structured way to report \textbf{abnormal situations} without immediately crashing a program.
Instead of returning special values like \texttt{-1} everywhere, Java can:
\begin{itemize}
  \item stop the normal flow,
  \item ``bubble up'' the problem to a place that can handle it,
  \item and still clean up resources reliably.
\end{itemize}
In this lecture you will learn what exceptions are, how to handle them, and how to create simple custom exceptions for clearer programs.
\notesources{\OOPSUnitIIISources}

\section{Prerequisite Bridge}
Before reading the new material in depth, do a fast recall of the older ideas listed below.
\begin{itemize}
  \item \textbf{Control flow}: try, catch, and finally are still flow-control ideas for abnormal situations. Main reference: \texttt{Unit 01 / Lecture 04 slides/notes}.
  \item \textbf{Validation design}: Encapsulated objects should reject invalid state changes in a disciplined way. Main reference: \texttt{Unit 02 / Lecture 01 slides/notes}.
\end{itemize}
This lecture only revisits those ideas briefly so that the main explanation can stay focused on the new concept sequence.
\notesources{\OOPSUnitIIISources}

\section{Key Terms (Quick)}
\begin{itemize}
  \item \textbf{Exception}: abnormal condition a program can potentially handle.
  \item \textbf{Error}: serious problem typically not recoverable (e.g., \texttt{OutOfMemoryError}).
  \item \textbf{Checked exception}: must be caught or declared (e.g., \texttt{IOException}).
  \item \textbf{Unchecked exception}: subclass of \texttt{RuntimeException} (e.g., \texttt{NullPointerException}).
  \item \textbf{Stack trace}: call history showing where an exception occurred.
  \item \textbf{Handler}: catch block that responds to an exception.
\end{itemize}

\section{Detailed Notes}
\subsection{Exceptions vs Errors}
\begin{itemize}
  \item \textbf{Exceptions} are conditions that can be handled: invalid input, missing file, etc.
  \item \textbf{Errors} are typically JVM-level problems: memory exhaustion, serious linkage errors, etc.
\end{itemize}

\paragraph{Rule of thumb.} Catch exceptions you can handle meaningfully. Do not catch \texttt{Error} in normal application code.

\subsection{What happens when an exception occurs? (stack unwinding)}
When an exception is thrown:
\begin{enumerate}
  \item Java creates an exception object (with message + stack trace).
  \item The current method stops executing immediately.
  \item Java looks for a matching \texttt{catch} in the current method.
  \item If not found, the exception goes to the caller (method that called it).
  \item This continues until a catch is found or the program terminates.
\end{enumerate}
This process is called \textbf{stack unwinding}.

\subsection{How to read a stack trace (very important skill)}
When you see a stack trace, do this:
\begin{enumerate}
  \item Find the first line: exception type + message.
  \item Look for the first stack frame that points to \textbf{your code} (your file/class).
  \item That line number is usually where the exception was thrown.
  \item If you see \texttt{Caused by:}, read the \textbf{root cause} first.
\end{enumerate}
Do not panic: stack traces are your friend.

\subsection{Checked vs unchecked exceptions (when to use which?)}
\small
\begin{tabular}{@{}p{0.24\linewidth}p{0.68\linewidth}@{}}
\toprule
\textbf{Category} & \textbf{Typical meaning} \\
\midrule
Checked (\texttt{Exception}) & {\raggedright Caller should handle it when recovery is realistic, e.g., \texttt{IOException}.} \\
Unchecked (\texttt{RuntimeException}) & {\raggedright Usually a programming mistake or invalid argument.\newline Runtime recovery is rare, e.g., \texttt{NullPointerException}.} \\
\bottomrule
\end{tabular}
\normalsize

\paragraph{Practical guideline.}
\begin{itemize}
  \item Use \textbf{checked} when the caller can take a meaningful alternate action (try another file, ask user again, retry).
  \item Use \textbf{unchecked} for invalid arguments and broken assumptions (caller passed wrong values).
\end{itemize}

\subsection{\texttt{try/catch/finally} and multiple catch}
Basic structure:
\begin{itemize}
  \item \texttt{try} contains code that may throw an exception.
  \item \texttt{catch} handles specific exception types.
  \item \texttt{finally} executes whether an exception occurs or not (useful for cleanup).
\end{itemize}

\textbf{Best practice:}
\begin{itemize}
  \item catch the most specific exception you can handle,
  \item avoid \texttt{catch (Exception e)} unless you rethrow/log properly,
  \item never silently swallow exceptions.
\end{itemize}

\paragraph{Catch order matters.}
Always catch the more specific exceptions first. If you catch \texttt{Exception} first, other catches become unreachable.

\paragraph{Multi-catch (one handler for multiple types).}
\begin{lstlisting}
try {
    // ...
} catch (NumberFormatException | IllegalArgumentException ex) {
    System.out.println("Bad input: " + ex.getMessage());
}
\end{lstlisting}
Use multi-catch when the handling is identical.

\subsection{Full example: \texttt{try/catch/finally} with multiple catches}
\textbf{Problem statement:} read two integers from the user and print \texttt{a/b}.
We want the program to:
\begin{itemize}
  \item handle invalid integer input (e.g., user types \texttt{abc}),
  \item handle divide-by-zero cleanly,
  \item always perform cleanup (close \texttt{Scanner}) and print a final message.
\end{itemize}

\subsubsection*{Full code (copy, compile, run)}
\begin{lstlisting}
import java.util.Scanner;

public class DivideDemo {
    public static void main(String[] args) {
        Scanner sc = new Scanner(System.in);

        try {
            System.out.print("Enter a: ");
            int a = Integer.parseInt(sc.nextLine());

            System.out.print("Enter b: ");
            int b = Integer.parseInt(sc.nextLine());

            int q = a / b; // may throw ArithmeticException
            System.out.println("a/b = " + q);

        } catch (NumberFormatException ex) {
            System.out.println("Please enter valid integers.");

        } catch (ArithmeticException ex) {
            System.out.println("Cannot divide by zero.");

        } finally {
            sc.close();
            System.out.println("Cleanup done (finally).");
        }
    }
}
\end{lstlisting}

\subsubsection*{Detailed explanation (what each block does)}
\textbf{1) \texttt{try} block:}
\begin{itemize}
  \item Reads input using \texttt{sc.nextLine()}.
  \item Converts to integer using \texttt{Integer.parseInt(...)}.
  \item Computes the quotient using \texttt{a / b}.
\end{itemize}

\textbf{What can go wrong inside \texttt{try}?}
\begin{itemize}
  \item If the user types \texttt{abc}, then \texttt{Integer.parseInt("abc")} throws \texttt{NumberFormatException}.
  \item If the user enters \texttt{b = 0}, then \texttt{a / 0} throws \texttt{ArithmeticException}.
\end{itemize}

\textbf{2) \texttt{catch} blocks:}
\begin{itemize}
  \item A catch block runs only if an exception of that type (or subtype) occurred in the try block.
  \item Java chooses the \textbf{first matching catch} from top to bottom.
\end{itemize}

\textbf{3) \texttt{finally} block:}
\begin{itemize}
  \item Runs \textbf{after} the try/catch finishes --- whether an exception happened or not.
  \item Used for \textbf{cleanup}: closing files, DB connections, scanners, releasing locks.
  \item It also runs even if you \texttt{return} inside \texttt{try} or \texttt{catch}.
\end{itemize}
\textbf{Rare exceptions:} \texttt{finally} may not run if the JVM exits abruptly (e.g., \texttt{System.exit(0)}) or the process is killed.

\subsubsection*{Sequence of execution (step-by-step)}
Students often memorize the syntax but get confused about \textbf{when} each block runs.
Use this mental model:
\begin{enumerate}
  \item Control enters the \texttt{try} block and executes statements \textbf{top to bottom}.
  \item If no exception occurs in the \texttt{try} block:
    \begin{itemize}
      \item the \texttt{catch} blocks are skipped,
      \item the \texttt{finally} block executes,
      \item execution continues after the whole try/catch/finally structure.
    \end{itemize}
  \item If an exception occurs at some line inside \texttt{try}:
    \begin{itemize}
      \item the remaining statements in \texttt{try} are skipped immediately,
      \item Java searches catch blocks from \textbf{top to bottom} and selects the \textbf{first matching} one,
      \item the selected \texttt{catch} runs (others are skipped),
      \item then \texttt{finally} runs,
      \item then execution continues after the structure.
    \end{itemize}
  \item If no \texttt{catch} matches the exception type:
    \begin{itemize}
      \item \texttt{finally} still runs,
      \item then the exception is \textbf{propagated} to the caller (the method exits abnormally).
    \end{itemize}
\end{enumerate}

\paragraph{Catch matching rule (very important).}
A catch matches if it has the same exception type or a \textbf{superclass} of the thrown exception.
That is why catch order matters: specific first, generic last.

\paragraph{Return + finally.}
If you \texttt{return} from \texttt{try} or \texttt{catch}, Java still executes \texttt{finally} before returning.

\subsubsection*{Sample runs (what you should observe)}
\textbf{Case A: valid input}
\begin{lstlisting}[language=]
Enter a: 10
Enter b: 2
a/b = 5
Cleanup done (finally).
\end{lstlisting}

\textbf{Case B: invalid number}
\begin{lstlisting}[language=]
Enter a: abc
Please enter valid integers.
Cleanup done (finally).
\end{lstlisting}

\textbf{Case C: divide by zero}
\begin{lstlisting}[language=]
Enter a: 10
Enter b: 0
Cannot divide by zero.
Cleanup done (finally).
\end{lstlisting}

\paragraph{Note about closing \texttt{Scanner}.}
Closing a \texttt{Scanner} created from \texttt{System.in} also closes the underlying input stream.
In a longer program that continues reading input later, you typically create one scanner and reuse it rather than repeatedly closing/reopening it.

\subsection{\texttt{throw} vs \texttt{throws}}
\begin{itemize}
  \item \texttt{throw} actually throws an exception now.
  \item \texttt{throws} declares that a method may throw an exception (caller must handle/declare).
\end{itemize}

\subsubsection*{Example A: \texttt{throw} (explicitly throwing now)}
\textbf{Goal:} stop the method immediately when an argument is invalid.
\begin{lstlisting}
static int divide(int a, int b) {
    if (b == 0) {
        throw new IllegalArgumentException("b must not be 0");
    }
    return a / b;
}
\end{lstlisting}
\textbf{What happens at runtime?}
\begin{itemize}
  \item If \texttt{b == 0}, Java creates an \texttt{IllegalArgumentException} object and throws it.
  \item The method exits immediately; the \texttt{return} line is skipped.
  \item Control jumps to the nearest matching \texttt{catch} (if any), otherwise the program terminates with a stack trace.
\end{itemize}
\textbf{When to use this pattern:} for invalid arguments and broken assumptions (\path|IllegalArgumentException| is common).

\subsubsection*{Example B: \texttt{throws} (declaring a checked exception)}
\textbf{Goal:} let the caller decide what to do if the operation fails.
\begin{lstlisting}
static void pauseOneSecond() throws InterruptedException {
    Thread.sleep(1000); // may throw InterruptedException (checked)
}
\end{lstlisting}
\textbf{Why \texttt{throws}?} \texttt{InterruptedException} is a \textbf{checked} exception, so the compiler forces you to either:
\begin{itemize}
  \item \textbf{catch} it, or
  \item \textbf{declare} it using \texttt{throws} in your own method signature.
\end{itemize}

\paragraph{Option 1: handle it (catch).}
\begin{lstlisting}
public static void main(String[] args) {
    try {
        pauseOneSecond();
        System.out.println("Done");
    } catch (InterruptedException ex) {
        System.out.println("Interrupted");
    }
}
\end{lstlisting}
\textbf{Meaning:} you handle the problem here (show message, cleanup, retry, etc.).

\paragraph{Option 2: propagate it (declare throws).}
\begin{lstlisting}
public static void main(String[] args) throws InterruptedException {
    pauseOneSecond();
    System.out.println("Done");
}
\end{lstlisting}
\textbf{Meaning:} you are not handling it here; you are pushing responsibility to your caller (for \texttt{main}, this means the JVM may print a stack trace if it happens).

\paragraph{Where should you use \texttt{throws}?}
Use \texttt{throws} when the method cannot handle the problem meaningfully and it is better to let a higher-level caller decide what to do.

\subsection{Custom exception design (basic, but practical)}
Create a custom exception when you want a clear domain-specific error message and type.
Good custom exceptions:
\begin{itemize}
  \item are named after the problem (\texttt{InvalidBookingException}),
  \item include a helpful message for debugging/user feedback,
  \item (optionally) wrap the original cause exception.
\end{itemize}

\paragraph{Checked or unchecked for custom exceptions?}
\begin{itemize}
  \item Use \textbf{unchecked} (\texttt{extends RuntimeException}) for validation and invalid arguments.
  \item Use \textbf{checked} (\texttt{extends Exception}) when callers should be forced to handle/recover.
\end{itemize}
\notesources{\OOPSUnitIIISources}

\section{Running Project Connection}
Use Exception Handling (Core) to make Campus Resource Manager safer, more modular, and easier to recover when something goes wrong.
\begin{itemize}
  \item Use Campus Resource Manager as the running case study, or map the same idea to your own project domain if you are using another example.
  \item Ask where today's idea should live: model, service, DAO, UI, or team workflow.
  \item Use the lecture example as a smaller version of the same design choice you will make in the capstone.
\end{itemize}
\notesources{\OOPSUnitIIISources}

\section{Common Mistakes and Confusions}
\begin{itemize}
  \item Using exceptions for normal control flow.
  \item Swallowing an exception silently.
  \item Showing raw technical exception details to end users.
\end{itemize}
\notesources{\OOPSUnitIIISources}

\section{Worked Example (tutorial): build a safe validation flow}
We will design a small method that validates input and reports problems clearly.

\subsection*{Step 1: Create a small validation exception (unchecked)}
\begin{lstlisting}
class ValidationException extends RuntimeException {
    public ValidationException(String message) {
        super(message);
    }

    public ValidationException(String message, Throwable cause) {
        super(message, cause);
    }
}
\end{lstlisting}
Why unchecked? If a user enters invalid input, we usually show a message and ask again. We don't want to clutter every method with mandatory \texttt{throws}.

\subsection*{Step 2: Write a parsing method that throws a clear exception}
\begin{lstlisting}
class Parsers {
    static int parsePositiveInt(String s) {
        if (s == null) throw new ValidationException("Input is required.");
        s = s.trim();
        if (s.isEmpty()) throw new ValidationException("Input cannot be empty.");

        try {
            int x = Integer.parseInt(s);
            if (x <= 0) throw new ValidationException("Number must be > 0.");
            return x;
        } catch (NumberFormatException ex) {
            throw new ValidationException("Not a valid integer: " + s, ex);
        }
    }
}
\end{lstlisting}
Notice:
\begin{itemize}
  \item We validate \textbf{null} and \textbf{empty} before parsing.
  \item We catch \texttt{NumberFormatException} and throw a better message.
  \item We keep the original cause (\texttt{ex}) for debugging.
\end{itemize}

\subsection*{Step 3: Use it in \texttt{main} and handle only at the boundary}
\begin{lstlisting}
import java.util.Scanner;

public class Demo {
    public static void main(String[] args) {
        Scanner sc = new Scanner(System.in);

        try {
            System.out.print("Enter number of tickets: ");
            int tickets = Parsers.parsePositiveInt(sc.nextLine());
            System.out.println("You booked " + tickets + " tickets.");
        } catch (ValidationException ex) {
            System.out.println("Input error: " + ex.getMessage());
        } finally {
            sc.close();
        }
    }
}
\end{lstlisting}
This is a clean pattern: validate/throw in low-level methods, and handle at the boundary (UI/CLI layer) where you can show messages.

\section{Demo / Observation Task}
Use this block as a short reproducible demo or observation task.
\begin{itemize}
  \item Reproduce the lecture's main example around \textit{Exceptions vs Errors} once without looking at the solution first.
  \item Change one input, rule, or design choice in Campus Resource Manager and predict the result before checking it.
  \item Record one success case, one edge case, and one failure case so the idea becomes observable instead of purely theoretical.
\end{itemize}
\notesources{\OOPSUnitIIISources}

\section{Worked Example: custom checked exception (recoverable case)}
\begin{lstlisting}
class InsufficientFundsException extends Exception {
    public InsufficientFundsException(String message) {
        super(message);
    }
}

class BankAccount {
    private double balance;

    public BankAccount(double openingBalance) {
        if (openingBalance < 0) throw new IllegalArgumentException("openingBalance");
        balance = openingBalance;
    }

    public void withdraw(double amount) throws InsufficientFundsException {
        if (amount <= 0) throw new IllegalArgumentException("amount");
        if (amount > balance) throw new InsufficientFundsException("Need " + amount + ", have " + balance);
        balance -= amount;
    }
}
\end{lstlisting}

\begin{lstlisting}
public class Demo {
    public static void main(String[] args) {
        BankAccount a = new BankAccount(100.0);
        try {
            a.withdraw(150.0);
        } catch (InsufficientFundsException e) {
            System.out.println("Withdraw failed: " + e.getMessage());
        }
    }
}
\end{lstlisting}
Here the caller can decide an alternate action: deposit more money, withdraw less, or show a message. That makes a checked exception reasonable.

\section{Demo / Observation Task}
Use this block as a short reproducible demo or observation task.
\begin{itemize}
  \item Reproduce the lecture's main example around \textit{Exceptions vs Errors} once without looking at the solution first.
  \item Change one input, rule, or design choice in Campus Resource Manager and predict the result before checking it.
  \item Record one success case, one edge case, and one failure case so the idea becomes observable instead of purely theoretical.
\end{itemize}
\notesources{\OOPSUnitIIISources}

\section{Practice Check (with brief solutions)}
\begin{enumerate}
  \item Name one checked and one unchecked exception.\par
  \textbf{Answer:} checked: \texttt{IOException}; unchecked: \texttt{NullPointerException}.

  \item What is the difference between \texttt{throw} and \texttt{throws}?\par
  \textbf{Answer:} \texttt{throw} throws now; \texttt{throws} declares possibility in method signature.

  \item When does \texttt{finally} execute?\par
  \textbf{Answer:} almost always (whether exception happens or not), used for cleanup.
\end{enumerate}

\section{Common bad practices (avoid)}
\begin{itemize}
  \item \textbf{Empty catch block}: hides bugs and makes debugging impossible.
  \item \textbf{Catching \texttt{Exception} everywhere}: can swallow unrelated problems.
  \item \textbf{Using exceptions for normal flow}: exceptions are slower and reduce clarity.
  \item \textbf{Throwing generic \texttt{Exception}}: prefer specific types with good messages.
\end{itemize}

\section{Self-Study Practice (no solutions)}
\begin{enumerate}
  \item Write a method \texttt{parsePositiveInt(String s)} that throws \texttt{IllegalArgumentException} on invalid input.
  \item Write a program that reads 2 integers and prints division result; handle divide-by-zero properly.
  \item Design a custom exception \texttt{InvalidBookingException} for a booking system.
\end{enumerate}

\section{Summary (what to remember)}
\begin{itemize}
  \item Exceptions separate \textbf{normal flow} from \textbf{error handling}.
  \item Checked exceptions force callers to handle recoverable problems.
  \item Unchecked exceptions are used for invalid arguments and programming errors.
  \item Handle exceptions at boundaries (UI/CLI); throw/propagate from deeper logic.
  \item Learn to read stack traces: find your code line and the root cause.
\end{itemize}

\section{Further Study}
\begin{itemize}
  \item Revisit the prerequisite references if any recall bullet still feels weak.
  \item Rework the lecture example with one changed assumption, input, or design choice.
  \item Read the textbook and documentation references listed in the source line for a second explanation of the same topic.
  \item Apply the same idea once inside Campus Resource Manager so that the concept moves from theory to design practice.
\end{itemize}
\notesources{\OOPSUnitIIISources}

\section{Exit Questions (with sample answers)}
\textbf{Q1.} Why should we catch specific exceptions instead of \texttt{Exception}?\par
\textbf{Sample answer:} Catching specific exceptions helps you handle known cases correctly and avoids hiding unrelated bugs. Catching \texttt{Exception} broadly can swallow errors and make debugging hard.

\vspace{0.6em}
\textbf{Q2.} What does checked exception mean?\par
\textbf{Sample answer:} A checked exception must be handled (caught) or declared with \texttt{throws}; otherwise the code will not compile.

\vspace{0.6em}
\textbf{Q3.} Why does Java separate checked and unchecked exceptions?\par
\textbf{Sample answer:} It distinguishes recoverable handled cases from bugs and invalid assumptions.

\end{document}
